% LLNCS LaTeX paper — Springer Proceedings format
% Methodology-only version (5-page target).
\documentclass[runningheads]{llncs}

\usepackage[T1]{fontenc}
\usepackage{graphicx}
\usepackage{booktabs}
\usepackage{amsmath}
\usepackage{hyperref}
\usepackage{subcaption}
\renewcommand\UrlFont{\color{blue}\rmfamily}
\usepackage{color}

\begin{document}

\title{Enhancing Customer Experience and Brand Perception: An NLP Pipeline for Hotel Review Analysis}
\titlerunning{Hotel Review NLP Pipeline: Methodology}

\author{
  Andr\'{e} Pires\inst{1}\and Manuel Sampaio\inst{1}\and Max Salomon\inst{1}\and
  Pedro Correia\inst{1}\and Pedro Meireles\inst{1}\and Mi\l{}osz Koz\l{}uk\inst{1}
}
\authorrunning{A. Pires et al.}
\institute{FEP, Faculty of Economics of the University of Porto, Portugal}

\maketitle

\begin{abstract}
This paper describes an NLP pipeline designed to analyse 999 Booking.com
guest reviews of a boutique hotel in Porto, Portugal, collected between
March 2023 and April 2026. The pipeline integrates exploratory analysis,
K-Means customer segmentation, supervised sentiment classification against
a lexicon baseline, Latent Dirichlet Allocation topic modelling, and
rule-based Aspect-Based Sentiment Analysis. For each stage we document the
design decisions, the alternatives considered, and the justifications
grounded in the properties of the data and the constraints of the
hospitality-review setting. Empirical results are reserved for the
accompanying video.
\keywords{Hotel reviews \and Sentiment analysis \and Topic modelling
\and Customer segmentation \and Aspect-Based Sentiment Analysis \and NLP}
\end{abstract}

\section{Introduction}

Online guest reviews have become a primary channel of customer
intelligence in hospitality. Mehraliyev et al.~\cite{mehraliyev2022}
note that most hospitality sentiment studies rely on off-the-shelf
lexicon tools or a single supervised classifier, without integrating
the complementary views offered by segmentation, topic modelling, and
aspect-level analysis. Our work designs a pipeline combining these
views, applied to one boutique property in Porto.

The contribution is deliberately methodological. For each stage we
record the choices made, the data properties that motivated them, and
the alternatives that were rejected. Empirical findings produced by
those choices are presented in the accompanying video. The paper is
organised as a chain of design justifications: exploratory analysis
informs the features used in segmentation; the observed class imbalance
informs the sentiment classifier configuration; the absence of aspect
annotations informs a rule-based approach to ABSA.

\section{Data Collection and Preparation}

Reviews were collected via a Selenium-based scraper targeting one hotel.
Scraping was chosen because Booking.com exposes no public API for
third-party review retrieval; Selenium is the standard workaround in prior
hospitality NLP work. The scraper handled dynamic pagination, modal
dialogs, and content-based deduplication. Captured fields included
reviewer country, room type, nights, traveller type, positive and
negative review text, and numerical score. Reviewer names were used only
for deduplication and otherwise excluded, following GDPR
data-minimisation.

Cleaning was kept minimal and transparent: date parsing to year-month
strings, extraction of the integer nights field, and concatenation of the
positive and negative fields into a unified review text. The
concatenation is deliberate: Booking.com's dual-field structure makes the
two fields semantically complementary rather than independent, so
separating them would double-count reviewers and distort downstream
frequency statistics.

Two data properties then constrained the rest of the pipeline. A
non-trivial fraction of reviews are not in English, and no labelled
multilingual hotel-review training set of sufficient size was
available; we therefore applied \texttt{langdetect} and Google Translate
to non-English reviews so that downstream classical models could operate
in a single language space. Booking.com also publishes no ground-truth
sentiment labels, only numerical scores, so we derived three-class
labels via conservative thresholds (positive: ${\geq}\,8$; neutral: 6--7;
negative: ${<}\,6$). The proxy is standard in hospitality NLP but
embeds a label-signal mismatch discussed in Sec.~\ref{sec:sentiment}.

\section{Exploratory Data Analysis}
\label{sec:eda}

Exploratory analysis here is a design step rather than a findings step:
its role is to surface data properties that motivate subsequent choices.
Four signals are consequential. The score distribution is heavily
left-skewed, implying severe class imbalance under any score-based
labelling scheme, which shapes the sentiment classifier configuration
(Sec.~\ref{sec:sentiment}). Traveller type partitions guests into groups
of markedly different size and behaviour, making it an informative
clustering feature (Sec.~\ref{sec:segmentation}) rather than a confound.
Word-frequency profiles for positive and negative subsets surface a
compact set of recurring concerns (room amenities, noise, connectivity,
staff, location, breakfast) that guided the eight aspects used in ABSA
(Sec.~\ref{sec:absa}). Review volume is unevenly distributed across 38
months, so models sensitive to temporal drift were avoided and
stratified cross-validation was adopted downstream.

\section{Customer Segmentation}
\label{sec:segmentation}

Segmentation is included because marketing actions (loyalty,
family-targeted amenities, language-specific outreach) are naturally
formulated per segment rather than per review. K-Means was used for three
reasons: interpretability, scaling to the corpus size, and compatibility
with a mixed numeric / one-hot feature representation, without the
density assumptions that would favour Gaussian mixture clustering. DBSCAN
and hierarchical clustering were considered and rejected. DBSCAN performs
poorly on features with disparate scales even after standardisation, and
hierarchical clustering scales less gracefully while offering no
interpretability advantage.

Five features combine numerical guest score (the outcome), review word
count (a proxy for engagement depth), number of nights (a proxy for stay
type), traveller type in one-hot form, and a three-level ordinal encoding
of the derived sentiment label. Each captures a distinct behavioural
axis, and each is available from the cleaned data without further
engineering. All features were standardised with z-scores before
clustering; standardisation is mandatory for K-Means on mixed-scale inputs
since otherwise higher-magnitude features dominate the Euclidean
distance. The number of clusters was selected by elbow analysis on the
inertia curve over $k \in \{2,\ldots,8\}$, taking the smallest $k$ beyond
which inertia reduction flattened.

Two limitations follow from the design. Clustering is at the review
level, not the customer level, because Booking.com reviews are not linked
by user identifier after scraping; a guest with multiple stays therefore
contributes multiple vectors. The one-hot encoding of traveller type also
means resulting clusters partly recover the a priori traveller-type
grouping, which is acknowledged rather than hidden in
Sec.~\ref{sec:discussion}.

\section{Sentiment Analysis}
\label{sec:sentiment}

Three decisions shape the sentiment stage: feature representation, baseline choice, and reformulation of the task.

TF-IDF vectors (max 8{,}000 terms, unigram and bigram, sublinear term
frequency) paired with linear classifiers (Logistic Regression, LinearSVC)
form the feature representation. With fewer than one thousand labelled
reviews, the sample is well below the regime where transformer fine-tuning
reliably outperforms classical linear baselines for sentiment
classification. Linear models on TF-IDF also preserve interpretability,
a property that matters when downstream users are hotel operators rather
than ML engineers.

VADER~\cite{vader2014}, a rule-based lexicon analyser requiring no
training data, serves as the baseline. It is the standard reference point
in hospitality sentiment work, isolating the contribution of supervised
learning: any supervised model that fails to outperform VADER in-domain
offers little value beyond the lexicon itself. All supervised classifiers
used \texttt{class\_weight=`balanced'} to correct for the imbalance
surfaced in Sec.~\ref{sec:eda}, and 5-fold stratified cross-validation.
Stratification is required: simple K-fold splits on imbalanced data
produce folds in which minority classes are sparsely or unevenly
represented.

The task was reformulated as binary. The three-class formulation
inherited from score-derived labels carries a structural problem: a
review scored 7 is linguistically close to a review scored 8, and the
neutral / positive boundary has no counterpart in the text, producing a
label-signal mismatch that caps achievable performance. Collapsing to
non-positive versus positive (threshold at score 8) eliminates the
troubled boundary and yields the primary deployable configuration.
Binary reformulation is common in hotel-review NLP for exactly this
reason.

Alongside the classical models, \texttt{xlm-roberta-base}~\cite{xlmroberta2020}
was fine-tuned for 2 epochs as a multilingual baseline, to verify that
the classical TF-IDF route remains competitive at this data scale and
document the path available if the dataset were later scaled across
properties. Monolingual BERT variants were rejected because they would
force discarding the non-English subset.

\section{Topic Modelling}

Topic modelling serves a distinct purpose from segmentation: segmentation
partitions \textit{reviewers}, whereas topic modelling partitions
\textit{themes}. Latent Dirichlet Allocation~\cite{blei2003}, the
foundational model for discovering interpretable topic-word distributions
over a text corpus, was chosen over BERTopic on two grounds. Its topic
distributions are directly readable by non-ML stakeholders, which fits
the methodology paper's communicative goal, and it does not require
pre-trained language-model embeddings whose multilingual behaviour on
the data would be difficult to audit.

Vocabulary was built with CountVectorizer using unigram and bigram
features, min\_df=3, max\_df=0.90, and max\_features=3{,}000. The min\_df
threshold removes idiosyncratic terms that would produce spurious
single-review topics; max\_df removes near-universal terms (`hotel',
`room', `stay') that provide no discriminative signal; max\_features
controls model complexity. A hospitality-domain stop-word list extended
the default English list with multilingual filler terms identified during
EDA. The number of topics was selected by comparing perplexity across
$k \in \{4, 6, 8\}$, with the smaller $k$ under-partitioning and the
larger $k$ fragmenting identifiable themes into near-duplicates.

\section{Aspect-Based Sentiment Analysis}
\label{sec:absa}

ABSA operates at sub-review granularity, attributing sentiment to
specific hotel features. Zhang et al.~\cite{zhang2022} identify a
spectrum of ABSA methods spanning rule and lexicon approaches,
supervised neural models, and large language model prompting. A
rule-based method was adopted for three reasons specific to our setting.
Supervised neural ABSA requires aspect-level annotations that we do not
have: the data contains numerical scores but no span or aspect-polarity
labels, and annotating roughly one thousand reviews across eight
aspects was out of scope. Pre-trained ABSA models from the SemEval
restaurant or laptop tracks do not transfer cleanly to hospitality
vocabulary. Rule-based extraction is also transparent: aspect set,
keyword lists, and scoring procedure can all be audited by a non-ML
reader, aligning with the pipeline's interpretability goal.

The eight aspects were not chosen a priori but derived from the
word-frequency analysis in Sec.~\ref{sec:eda}: the most recurrent
hotel-domain nouns and noun phrases were grouped into Location, Staff \&
Service, Room, Cleanliness, Breakfast, Noise, Value, and WiFi \&
Check-in. For each aspect a keyword list was built; each review text
was then scanned for keyword occurrences, and a context window of
${\pm}\,12$ tokens around each hit was scored with VADER. The window is a
standard setting in lexicon-based ABSA: small enough to exclude
unrelated sentiment from distant clauses, large enough to capture a
modifier-head relationship that may span a few tokens. Aspect mentions
were deduplicated per review so that a guest mentioning WiFi three
times contributes one instance per aspect. The output is, for each
review and each aspect, either a signed sentiment score or a null.

\section{Methodological Discussion}
\label{sec:discussion}

Several choices involve trade-offs worth making explicit. Linear models,
rule-based ABSA, and classical LDA were adopted despite stronger neural
alternatives, since the data size does not support reliable transformer
fine-tuning and since pipeline users are hotel operators rather than ML
practitioners; a richer dataset across properties would justify
revisiting the trade-off. Score-derived sentiment labels are a proxy,
embedding a label-signal mismatch that binary reformulation only
partially resolves; human-annotated labels on a subset would quantify
the gap directly. K-Means clusters at the review level because
Booking.com reviews carry no stable user identifier after scraping,
biasing centroids toward frequent reviewers. All findings and design
validations are conditional on one property; the pipeline runs unchanged
on other properties, and cross-property comparison is the most
informative next application. Booking.com also exposes no helpful-vote
counts or verified-purchase flags, precluding review helpfulness
prediction and fake-review detection tasks that are standard in related
TripAdvisor-based work~\cite{ott2011}.

{\footnotesize
\begin{thebibliography}{8}
\setlength{\itemsep}{0pt}
\sloppy

\bibitem{scikit-learn}
Pedregosa, F., et al.: Scikit-learn: Machine learning in Python.
\textit{JMLR} \textbf{12}, 2825--2830 (2011)

\bibitem{vader2014}
Hutto, C.J., Gilbert, E.: VADER: A parsimonious rule-based model for
sentiment analysis of social media text. In: \textit{Proc. ICWSM 2014}
(2014)

\bibitem{blei2003}
Blei, D.M., Ng, A.Y., Jordan, M.I.: Latent Dirichlet Allocation.
\textit{JMLR} \textbf{3}, 993--1022 (2003)

\bibitem{xlmroberta2020}
Conneau, A., et al.: Unsupervised cross-lingual representation learning
at scale. In: \textit{Proc. ACL 2020}, pp. 8440--8451 (2020)

\bibitem{mehraliyev2022}
Mehraliyev, F., Chan, I.C.C., Kirilenko, A.P.: Sentiment analysis in
hospitality and tourism: a thematic and methodological review.
\textit{Int. J. Contemp. Hosp. Manag.} \textbf{34}(1), 46--77 (2022)

\bibitem{zhang2022}
Zhang, W., Li, X., Deng, Y., Bing, L., Lam, W.: A survey on aspect-based
sentiment analysis: tasks, methods, and challenges. \textit{IEEE Trans.
Knowl. Data Eng.} \textbf{35}(11), 11019--11038 (2022)

\bibitem{ott2011}
Ott, M., Choi, Y., Cardie, C., Hancock, J.T.: Finding deceptive opinion
spam by any stretch of the imagination. In: \textit{Proc. ACL 2011},
pp. 309--319 (2011)

\end{thebibliography}
}

\end{document}
